% Blueprint roadmap for Garnett, Bounded Analytic Functions, Chapter II,
% Section 5. The source PDF is intentionally not tracked by git.
%
% The statements below are formalization-oriented paraphrases, not verbatim
% copies from the book.

\chapter{Canonical factorization in the Nevanlinna class}

\section{Chapter I prerequisites}

\begin{theorem}[Garnett I.5.3: Fatou theorem, $p=1$ disc form]
  \label{thm:I_5_3}
  \notready
  \uses{lem:I_5_4}
  Source: Garnett, Chapter I, Section 5, PDF pages 40--41.
  Let $\mu$ be a finite real or complex measure on
  $\partial\mathbb{D}$, with Lebesgue decomposition
  \[
    d\mu(\theta)=f(\theta)\,\frac{d\theta}{2\pi}+d\nu(\theta),
    \qquad \nu\perp d\theta.
  \]
  Then the Poisson integral of $\mu$ has nontangential boundary values equal
  to the density of the absolutely continuous part:
  \[
    \operatorname{nt}\!\lim_{z\to e^{i\theta}}P[\mu](z)=f(\theta)
  \]
  for Lebesgue almost every $\theta$. Garnett states the stronger theorem on
  the upper half-plane for all $1\le p\le\infty$; this is the $p=1$ consequence
  used later, transported to the disc by a Cayley transform.
\end{theorem}

\begin{theorem}[Harmonic majorants of subharmonic functions, Garnett I.6.7]
  \label{thm:I_6_7}
  \notready
  Source: Garnett, Chapter I, Section 6, PDF page 36.
  Let $v$ be subharmonic on $\mathbb{D}$. Then $v$ has a harmonic majorant on
  $\mathbb{D}$ if and only if
  \[
    \sup_{0<r<1}\frac{1}{2\pi}
      \int_0^{2\pi}v(r e^{i\theta})\,d\theta
    = \sup_{0<r<1} v_r(0) < \infty,
  \]
  where $v_r$ is the Poisson extension to $|z|<r$ of the boundary values of
  $v$ on $|z|=r$. When this condition holds, the least harmonic majorant $u$
  of $v$ is given by
  \[
    u(z)=\lim_{r\to 1}
      \frac{1}{2\pi}\int_0^{2\pi}
        P_{z/r}(\theta)v(r e^{i\theta})\,d\theta
      =\lim_{r\to 1}v_r(z).
  \]
\end{theorem}

\begin{lemma}[Garnett I.5.4: singular Poisson integrals vanish at the boundary]
  \label{lem:I_5_4}
  \notready
  Source: Garnett, Chapter I, Section 5, PDF pages 41--42.
  If $\nu$ is a finite measure on $\mathbb{R}$ singular with respect to
  Lebesgue measure, then its Poisson integral has nontangential boundary limit
  zero for Lebesgue almost every boundary point. This is the step that removes
  the singular part in the $p=1$ boundary-limit assertion of
  Theorem~\ref{thm:I_5_3}.
\end{lemma}

\section{Chapter II, Section 2: Blaschke products}

\begin{theorem}[Garnett II.2.1: harmonic majorants imply the Blaschke condition]
  \label{thm:II_2_1}
  \notready
  \uses{thm:I_6_7}
  Source: Garnett, Chapter II, Section 2, PDF page 63.
  Let $f$ be a nonzero analytic function on $\mathbb{D}$, and let
  $\{a_n\}$ be its zeros, counted with multiplicity. If $\log |f|$ has a
  harmonic majorant on $\mathbb{D}$, then
  \[
    \sum_n (1-|a_n|)<\infty.
  \]
  Moreover, if $f(0)\ne 0$ and $u$ is the least harmonic majorant of
  $\log |f|$, then
  \[
    \sum_n (1-|a_n|) \le u(0)-\log |f(0)|.
  \]
  By Theorem~\ref{thm:II_5_1}, the harmonic-majorant hypothesis is equivalent
  to $f$ belonging to the Nevanlinna class $N$.
  \begin{proof}
    After removing a possible zero at the origin, fix $r<1$ avoiding the
    moduli of the zeros and divide $f(rz)$ by the finite Blaschke product with
    the same zeros. The logarithm of the modulus of the resulting zero-free
    analytic function is harmonic, so the mean-value formula gives the
    familiar Jensen identity
    \[
      \log |f(0)|+\sum_{|a_j|<r}\log\frac{r}{|a_j|}
      =\frac{1}{2\pi}\int_0^{2\pi}\log|f(r e^{i\theta})|\,d\theta.
    \]
    Theorem~\ref{thm:I_6_7} identifies the limit of the right-hand side with
    $u(0)$. Letting $r\to1$ and using
    $1-|a_j|\le\log(1/|a_j|)$ proves both conclusions.
  \end{proof}
\end{theorem}

\begin{theorem}[Garnett II.2.2: convergence and boundary size of Blaschke products]
  \label{thm:II_2_2}
  \notready
  Source: Garnett, Chapter II, Section 2, PDF page 64.
  If $\{a_n\}\subset\mathbb{D}$ satisfies the Blaschke condition, then the
  associated Blaschke product converges uniformly on compact subsets of
  $\mathbb{D}$ to a bounded analytic function whose zeros are exactly the
  $a_n$ with the prescribed multiplicities, and whose boundary modulus is one
  almost everywhere.
\end{theorem}

\begin{theorem}[Garnett II.2.4: division by the Blaschke factor]
  \label{thm:II_2_4}
  \notready
  \uses{thm:I_6_7}
  Source: Garnett, Chapter II, Section 2, PDF page 66.
  Let $B$ be the Blaschke product built from part or all of the zero sequence
  of an analytic function $f$. Dividing by $B$ removes those zeros without
  increasing the relevant Hardy or boundary-size quantity; in particular the
  quotient remains in the same Hardy class in the applications used below.
\end{theorem}

\begin{definition}[Lean Blaschke factors]
  \label{def:blaschke_factor}
  \lean{BlaschkeFactor, normedBlaschkeFactor}
  \leanok
  \uses{thm:II_2_1, thm:II_2_2, thm:II_2_4}
  The current Lean project has the one-zero Blaschke factors. The later
  formalization target is to define infinite Blaschke products, prove compact
  convergence, prove the a.e. boundary-modulus theorem, and prove the quotient
  theorem needed for Hardy and Nevanlinna functions.
\end{definition}

\section{Chapter II, Section 4: boundary uniqueness and outer functions}

\begin{corollary}[Garnett II.4.2: boundary uniqueness in $H^p$]
  \label{cor:II_4_2}
  \notready
  Source: Garnett, Chapter II, Section 4, PDF page 75.
  Let $0<p\le\infty$ and let $f\in H^p(\mathbb{D})$. If the boundary function
  of $f$ vanishes on a subset of $\partial\mathbb{D}$ of positive Lebesgue
  measure, then $f$ is identically zero.
\end{corollary}

\begin{theorem}[Garnett II.4.4: construction from logarithmic data]
  \label{thm:II_4_4}
  \notready
  Source: Garnett, Chapter II, Section 4, PDF page 75.
  Given a nonnegative boundary function $\varphi$ with
  $\log\varphi\in L^1(\partial\mathbb{D})$, there is an outer analytic
  function $F$ on $\mathbb{D}$ whose a.e. boundary modulus is $\varphi$.
  Equivalently, $F$ is obtained by exponentiating the analytic function whose
  real part is the Poisson integral of $\log\varphi$.
\end{theorem}

\begin{theorem}[Outer uniqueness]
  \label{thm:outer_uniqueness}
  \notready
  If two outer functions have the same boundary modulus almost everywhere,
  then their quotient is a unimodular constant. Equivalently, the boundary
  modulus determines the outer factor up to the expected constant.
\end{theorem}

\begin{definition}[Outer factor]
  \label{def:outer_factor}
  \notready
  \uses{thm:II_4_4, thm:outer_uniqueness}
  The outer factor $F$ in the Nevanlinna factorization is the zero-free
  analytic factor built from the absolutely continuous part of the boundary
  measure representing the least harmonic majorant.
\end{definition}

\section{Chapter II, Section 5: Nevanlinna class}

\begin{theorem}[Garnett II.5.1: harmonic-majorant characterization of $N$]
  \label{thm:II_5_1}
  \notready
  \uses{thm:log_abs_subharmonic, thm:I_6_7,
    thm:positive_harmonic_poisson_measure}
  Source: Garnett, Chapter II, Section 5, PDF page 78.
  For an analytic function $f$ on $\mathbb{D}$, the following are equivalent:
  $f$ belongs to the Nevanlinna class $N$, and the subharmonic function
  $\log |f|$ admits a harmonic majorant on $\mathbb{D}$. When such a majorant
  exists, there is a least harmonic majorant.
\end{theorem}

\begin{theorem}[Logarithm of modulus is subharmonic]
  \label{thm:log_abs_subharmonic}
  \lean{logNormBot_comp_analytic_subharmonicOn_scalar,
    logNormBot_comp_analytic_subharmonicOn_banach}
  \leanok
  If $f$ is analytic on an open subset of $\mathbb{C}$, then
  $z\mapsto \log \lVert f(z)\rVert$, interpreted as $-\infty$ at zeros, is
  subharmonic there. This has been formalized both for complex-valued
  functions and, more generally, for functions taking values in a complex
  Banach space.
  \begin{proof}
    \proves{thm:log_abs_subharmonic}
    \leanok
    The Lean proofs are complete in the declarations linked above.
  \end{proof}
\end{theorem}

\begin{theorem}[Garnett I.3.5(c): positive harmonic functions as Poisson integrals, disc form]
  \label{thm:positive_harmonic_poisson_measure}
  \notready
  Source: Garnett, Chapter I, Section 3, Theorem 3.5(c), especially the
  disc-version argument in its proof, printed pages 17--19 (PDF pages
  approximately 29--31).
  If $u$ is a positive harmonic function on $\mathbb{D}$, then there is a
  finite positive measure $\mu$ on $\partial\mathbb{D}$ such that
  $u=P[\mu]$. Conversely, the Poisson integral of any finite positive boundary
  measure is positive and harmonic.
\end{theorem}

\begin{lemma}[Garnett II.5.2: remove the zeros of a Nevanlinna function]
  \label{lem:II_5_2}
  \notready
  \uses{thm:II_5_1, thm:II_2_1, thm:II_2_2, thm:II_2_4}
  Source: Garnett, Chapter II, Section 5, PDF page 78.
  Let $f\in N$ be nonzero and let $B$ be the Blaschke product formed from the
  zeros of $f$, counted with multiplicity. Then $B$ converges, the quotient
  $g=f/B$ is analytic and zero-free on $\mathbb{D}$, and $g$ still belongs to
  the Nevanlinna class. Moreover, $\log|g|$ is the least harmonic majorant of
  $\log|f|$ on $\mathbb{D}$.
\end{lemma}

\begin{corollary}[Garnett II.3.2: $H^\infty$ boundary values, disc form]
  \label{cor:II_3_2}
  \notready
  Source: Garnett, Chapter II, Section 3, PDF page 67.
  Every bounded analytic function on $\mathbb{D}$ has nontangential boundary
  limits almost everywhere, and the resulting boundary function has essential
  supremum bounded by the analytic supremum norm. This is the forward
  direction of the $p=\infty$ case of Corollary II.3.2, transported from the
  upper half-plane to the disc by a Cayley transform.
\end{corollary}

\begin{theorem}[Garnett II.5.3: boundary values and the least-majorant measure]
  \label{thm:II_5_3}
  \notready
  \uses{thm:II_5_1, lem:II_5_2, thm:I_5_3,
    cor:II_3_2, cor:II_4_2,
    thm:positive_harmonic_poisson_measure}
  Source: Garnett, Chapter II, Section 5, PDF page 79.
  If $f\in N$ and $f$ is not identically zero, then $f$ has nontangential
  boundary values $f(e^{i\theta})$ almost everywhere, and
  \[
    \log|f(e^{i\theta})|\in L^1(d\theta).
  \]
  If $u$ is the least harmonic majorant of $\log|f|$, then
  \[
    u(z)=\int_{0}^{2\pi}P_z(\theta)\,d\mu(\theta),
  \]
  where the finite real representing measure has the decomposition
  \[
    d\mu(\theta)
      =\log|f(e^{i\theta})|\,\frac{d\theta}{2\pi}+d\mu_s(\theta),
  \]
  with $\mu_s$ singular with respect to Lebesgue measure.
\end{theorem}

\section{Singular factors}

\begin{theorem}[Jordan decomposition of singular measures]
  \label{thm:jordan_decomposition_singular_measure}
  \notready
  If $\nu$ is a finite real signed measure on $\partial\mathbb{D}$ and is
  singular with respect to Lebesgue measure, then its Jordan decomposition
  $\nu=\nu^+-\nu^-$ has $\nu^+$ and $\nu^-$ finite positive singular measures.
\end{theorem}

\begin{definition}[Singular function formula]
  \label{def:singular_function_formula}
  \notready
  For a finite positive singular measure $\sigma$ on $\partial\mathbb{D}$, the
  associated singular function is the analytic exponential
  \[
    S_\sigma(z)=\exp\left(-\int_{\partial\mathbb{D}}
      \frac{\zeta+z}{\zeta-z}\,d\sigma(\zeta)\right).
  \]
  It is bounded by one in the disc and has boundary modulus one almost
  everywhere.
\end{definition}

\begin{definition}[Singular factors]
  \label{def:singular_factors}
  \notready
  \uses{thm:jordan_decomposition_singular_measure,
    def:singular_function_formula, lem:I_5_4}
  If the singular part of the representing measure is
  $\nu_s=\nu_s^+-\nu_s^-$, then the canonical factorization uses the two
  singular functions $S_1=S_{\nu_s^+}$ and $S_2=S_{\nu_s^-}$.
\end{definition}

\section{Final factorization and uniqueness}

\begin{theorem}[Zeros determine the Blaschke factor]
  \label{thm:zeros_determine_blaschke}
  \notready
  The multiset of zeros in $\mathbb{D}$ determines the associated Blaschke
  product up to the chosen unimodular normalization. With a fixed convention
  for elementary factors, the product is uniquely determined.
\end{theorem}

\begin{theorem}[Boundary modulus determines the outer factor]
  \label{thm:boundary_modulus_determines_outer}
  \notready
  \uses{thm:outer_uniqueness}
  If two outer functions have the same a.e. boundary modulus, then they differ
  by multiplication by a constant of modulus one.
\end{theorem}

\begin{theorem}[Singular measure determines singular factors]
  \label{thm:singular_measure_determines_factors}
  \notready
  \uses{def:singular_factors}
  The positive and negative singular parts of the representing measure
  determine $S_1$ and $S_2$. Conversely, equality of the corresponding
  singular functions forces equality of the singular measures, modulo the
  normalization convention.
\end{theorem}

\begin{theorem}[Uniqueness of Nevanlinna factorization]
  \label{thm:nevanlinna_factorization_uniqueness}
  \notready
  \uses{thm:zeros_determine_blaschke,
    thm:boundary_modulus_determines_outer,
    thm:singular_measure_determines_factors}
  In the canonical factorization, the Blaschke factor is determined by the
  zero multiset, the outer factor by the a.e. boundary modulus, and the
  singular factors by the singular part of the representing measure. The only
  remaining ambiguity is the conventional unimodular constant.
\end{theorem}

\begin{theorem}[Garnett II.5.5: canonical factorization in $N$]
  \label{thm:nevanlinna_canonical_factorization}
  \notready
  \uses{thm:II_5_1, lem:II_5_2, thm:II_5_3, def:outer_factor,
    def:singular_factors, thm:nevanlinna_factorization_uniqueness}
  Source: Garnett, Chapter II, Section 5, PDF page 83.
  If $f\in N$ and $f$ is not identically zero, then there exist a unimodular
  constant $C$, a Blaschke product $B$, an outer function $F$, and singular
  functions $S_1,S_2$ such that
  \[
    f = C\,B\,F\,S_1/S_2.
  \]
  Conversely, every expression of this form belongs to the Nevanlinna class.
  The factors are unique after fixing the standard normalizations.
\end{theorem}

\begin{proof}
  \notready
  Remove the zeros using the Blaschke product. Apply the boundary-value and
  logarithmic-integrability theorem to the zero-free quotient. Split the
  representing measure of the least harmonic majorant into absolutely
  continuous and singular parts. The absolutely continuous part gives the
  outer factor, while the positive and negative singular parts give
  $S_1$ and $S_2$. The uniqueness statement follows by comparing zero sets,
  boundary moduli, and singular measures.
\end{proof}
